\section{Step 1 Identify the current sample in the supplied input}\label{fir-s01}

Read the input as seven chronological history samples followed by the new samples. \texttt{a0} addresses the oldest history value. At the first output, array position seven is current; at the next output, the base slides one halfword and the current value is original position eight. Coefficients stay in oldest-first order, so the last stored coefficient multiplies the current sample.

\subsection{Expected observation}
One output requires eight input halfwords, two output bytes and the same sixteen coefficient bytes used by every output.

\Needspace{11\baselineskip}
\subsection{Check your understanding}
Which array position is current for output zero?

\begin{itemize}
  \item \textbf{A} Zero.
  \item \textbf{B} Seven.
  \item \textbf{C} Eight.
\end{itemize}
